\documentclass[11pt, a4paper]{article}
\usepackage[utf8]{inputenc}
\usepackage{amsmath, amssymb}
\usepackage{geometry}
\geometry{margin=1in}
\usepackage{hyperref}
\usepackage{graphicx}
\usepackage{bm}

\title{\textbf{Fibras: Synthetic STED Fiber Field Learning Pipeline}}
\author{Pipeline Documentation}
\date{}

\begin{document}

\maketitle

\begin{abstract}
This document describes the current Fibras project as implemented in the codebase.
The project no longer centers on direct topological graph analysis of thresholded
volumes. Its main path is a supervised learning pipeline for 2D STED-like fiber
images: synthetic 3D fiber scenes are rendered through an empirical optical model,
projected into 2D microscopy slices, converted into dense distance, orientation,
and visibility targets, and used to train a residual U-Net. At inference time the
network predicts these fields on full-resolution images and a streamline tracker
converts the fields into a binary skeleton.
\end{abstract}

\section{Current Project Scope}

The active pipeline is organized around four tasks:
\begin{enumerate}
    \item Generate synthetic 2D STED training data from 3D fiber geometry.
    \item Optionally calibrate synthesis parameters from real STED image statistics.
    \item Train a 2D residual U-Net to predict fiber localization, orientation, and visibility fields.
    \item Run tiled inference on real or synthetic images and trace predicted fields into skeleton masks.
\end{enumerate}

The main scripts are:
\begin{itemize}
    \item \texttt{analyze\_real\_sted.py}: measure real-image intensity, foreground, skeleton, noise, stripe, and vignette statistics.
    \item \texttt{generate\_dataset.py}: build train/validation/test \texttt{.pt} datasets.
    \item \texttt{train.py}: train or evaluate \texttt{STEDResUNet2D}.
    \item \texttt{inference.py}: run single-image, batch, or decoder-calibration inference.
    \item \texttt{visual.py}: inspect datasets, predictions, and STED synthesis debug views.
\end{itemize}

\section{Data Representation}

Each generated sample stores an input image and a four-channel target tensor:
\begin{equation}
    Y = \left[D,\; C,\; S,\; V\right],
\end{equation}
where $D$ is a normalized Euclidean-distance-transform-like centerline confidence,
$(C,S)$ are double-angle orientation channels, and $V$ is an axial visibility map.
For 2D samples the tensor layout on disk is:
\begin{equation}
    X \in \mathbb{R}^{1 \times 1 \times H \times W},
    \qquad
    Y \in \mathbb{R}^{4 \times 1 \times H \times W}.
\end{equation}

The orientation channels encode an unoriented tangent: a fiber direction
$\bm{u} = (u_x, u_y)$ and its opposite $-\bm{u}$ must be equivalent. The code
therefore maps vectors to double-angle coordinates:
\begin{equation}
    C = \cos(2\theta) = \frac{u_x^2 - u_y^2}{u_x^2 + u_y^2},
    \qquad
    S = \sin(2\theta) = \frac{2u_xu_y}{u_x^2 + u_y^2}.
\end{equation}
This removes the polarity ambiguity that would otherwise make tangent regression
unstable.

\section{Synthetic Geometry}

The geometry layer is built from \texttt{FiberSegment} objects, each holding a
start point, end point, and thickness multiplier. The current generator in
\texttt{src/synthesis.py} is:
\begin{itemize}
    \item \texttt{RandomWalkGenerator}: constrained short fibers with bounded turn angles.
\end{itemize}

The current 2D STED path primarily uses constrained 3D random walks. A synthetic
fiber is generated in a volume $(H,W,Z)$, but only its optical projection into a
selected focal slice is supervised. The generator samples step length, number of
steps, initial direction, turn angle, and bundle size. Bundles are formed by
adding jittered optical copies of the same core geometry:
\begin{equation}
    \tilde{\bm{p}} = \bm{p} + \bm{\epsilon}, \qquad
    \bm{\epsilon} \sim \mathcal{N}(0, \sigma_{\text{jitter}}^2 I).
\end{equation}
The unjittered core segments define labels; jittered copies affect only the
rendered optical appearance.

\section{Empirical STED Rasterization}

Rendering is handled by \texttt{EmpiricalRasterizer}. Segment density is computed
from the shortest distance to each segment and combined by maximum response for
basic rasterization:
\begin{equation}
    R(\bm{x}) =
    \max_i \exp\left(
        -\frac{d(\bm{x}, s_i)^2}{2\sigma_i^2}
    \right),
    \qquad
    \sigma_i = \sigma_{\text{base}} t_i.
\end{equation}

For STED-like 2D samples the code first rasterizes a 3D signal volume, then
collapses it to a focal slice with an axial defocus model. For a plane at axial
offset $\Delta z$ from the focal slice:
\begin{equation}
    w_z = \frac{1}{1 + \left(\frac{|\Delta z|}{d_{\text{of}}}\right)^2},
    \qquad
    \sigma_{\perp}(z) =
    \sqrt{\sigma_{\text{focus}}^2 +
    \left(\alpha_{\text{defocus}}|\Delta z|\right)^2}.
\end{equation}
Each plane is laterally blurred according to $\sigma_{\perp}(z)$, weighted by
$w_z$, and averaged into the final slice. The renderer also adds STED-specific
appearance effects: stochastic emission gaps, punctate modulation, monomer haze,
debris, striping, vignetting, Poisson photon noise, read noise, and a small DC
offset.

\section{Real-Data Calibration}

The calibrated synthesis mode is driven by \texttt{src/sted\_calibration.py}.
Real TIFF images are normalized, split into patches, and summarized with robust
foreground and artifact metrics:
\begin{itemize}
    \item intensity quantiles such as $p_{50}$, $p_{90}$, $p_{99}$, and $p_{99.9}$;
    \item foreground and skeleton fractions;
    \item component counts and approximate fiber widths;
    \item high-frequency background noise;
    \item low-frequency variation, vignetting, and stripe strength;
    \item orientation anisotropy.
\end{itemize}

These rows are stored as a calibration profile containing quantile summaries
globally and by real-data groups such as condition or DIV. During dataset
generation, \texttt{CalibrationSampler} samples a profile and converts measured
statistics into scene parameters: target skeleton density, base fiber width,
noise level, debris count, bundle probabilities, dynamic range, monomer haze
regime, stripe strength, and vignette edge range. The synthetic image is then
contrast-matched back to sampled real intensity percentiles.

\section{Target Construction}

\subsection{Distance and Orientation Targets}

\texttt{TargetFieldGenerator} computes dense supervision from vector segments.
For each pixel $\bm{x}$ and segment $s_i$, the closest point on the segment is
found and the minimum distance is retained:
\begin{equation}
    d(\bm{x}) = \min_i d(\bm{x}, s_i).
\end{equation}
The stored distance target is clipped and inverted:
\begin{equation}
    D(\bm{x}) = \max\left(0, 1 - \frac{d(\bm{x})}{d_{\max}}\right).
\end{equation}
The vector target at $\bm{x}$ is the tangent direction of the closest segment and
is converted to double-angle orientation channels for 2D training.

\subsection{Focus, Annotation, and Visibility Bands}

The 2D STED labels distinguish localization from visibility. A narrow focal slab
defines the centerlines the network should localize sharply. A broader annotation
band includes fibers that are optically visible but outside the strict focus
slab. The final EDT target blends these two maps:
\begin{equation}
    D_{\text{target}} =
    \max\left(D_{\text{focus}},\;
    \alpha_{\text{soft}} D_{\text{annotation}}\right).
\end{equation}
When the soft annotation exceeds the focus target, the corresponding orientation
target is taken from the annotation map.

The visibility target uses the same axial response as the renderer. A segment is
clipped to the axial band where its average optical weight remains above a floor.
Its contribution is rendered as a Gaussian tube weighted by mean axial visibility:
\begin{equation}
    V(\bm{x}) = \max_i
    \left[
    \bar{w}_{z,i}
    \exp\left(-\frac{d(\bm{x}, s_i)^2}{2\sigma_i^2}\right)
    \right].
\end{equation}
This target lets training down-weight ambiguous out-of-focus regions without
discarding them entirely.

\section{Network Architecture}

The active model is \texttt{STEDResUNet2D} in \texttt{src/model.py}. It is a
residual 2D U-Net with four encoder stages, a residual bottleneck, atrous spatial
pyramid pooling, and four decoder stages. The output is assembled from three
heads:
\begin{equation}
    f_{\theta}(X) =
    \left[\hat{D},\; \hat{C},\; \hat{S},\; \hat{L}_V\right],
\end{equation}
where $\hat{L}_V$ is a visibility logit. Training, evaluation, and inference are
explicitly 2D and use \texttt{STEDResUNet2D}.

\section{Training Objective}

\texttt{StedFieldLoss2D} combines four terms:
\begin{equation}
    \mathcal{L} =
    \mathcal{L}_{D}
    + \lambda_O \mathcal{L}_{O}
    + \lambda_V \mathcal{L}_{V}
    + \lambda_C(t) \mathcal{L}_{C}.
\end{equation}

The distance loss is visibility-weighted mean squared error:
\begin{equation}
    \mathcal{L}_{D} =
    \frac{\sum_{\bm{x}} c_V(\bm{x})
    \left(\hat{D}(\bm{x}) - D(\bm{x})\right)^2}
    {\sum_{\bm{x}} c_V(\bm{x}) + \epsilon},
    \qquad
    c_V = \max(V, V_{\min}).
\end{equation}

The orientation loss normalizes predicted and target double-angle vectors and
uses a confidence mask from the distance and visibility targets:
\begin{equation}
    \mathcal{L}_{O} =
    \frac{\sum_{\bm{x}} M(\bm{x})
    \left(1 - \hat{\bm{o}}(\bm{x}) \cdot \bm{o}(\bm{x})\right)}
    {\sum_{\bm{x}} M(\bm{x}) + \epsilon},
    \qquad
    M = \mathbb{1}[D > \tau_D] D c_V.
\end{equation}

Visibility is trained with binary cross entropy on logits. The centerline term is
a soft Dice-style error on high-power EDT ridges:
\begin{equation}
    \mathcal{L}_{C} =
    1 -
    \frac{2\sum_{\bm{x}} \hat{D}(\bm{x})^\gamma D(\bm{x})^\gamma + \epsilon}
    {\sum_{\bm{x}} \hat{D}(\bm{x})^\gamma +
     \sum_{\bm{x}} D(\bm{x})^\gamma + \epsilon}.
\end{equation}
The centerline weight is warmed up over the first training epochs, while
validation uses a fixed score so runs remain comparable.

\section{Inference and Streamline Extraction}

\texttt{inference.py} loads a \texttt{STEDResUNet2D} checkpoint, normalizes a
\texttt{.tif}, \texttt{.npy}, or \texttt{.pt} input image, and runs tiled
prediction with overlap blending. Tiling keeps full-resolution inference
practical for large microscopy images:
\begin{equation}
    \hat{Y} =
    \frac{\sum_k W_k \odot f_{\theta}(X_k)}
         {\sum_k W_k + \epsilon},
\end{equation}
where $W_k$ is a ramped blend window for tile $k$.

Predicted orientation channels are converted back to one arbitrary tangent
polarity:
\begin{equation}
    \hat{\theta} = \frac{1}{2}\operatorname{atan2}(\hat{S}, \hat{C}),
    \qquad
    \hat{\bm{u}} = (\cos\hat{\theta}, \sin\hat{\theta}).
\end{equation}
The tracker multiplies EDT by clipped visibility, skeletonizes the high-confidence
mask to create seeds, and integrates streamlines forward and backward through the
predicted vector field. Local vector polarity is flipped when necessary to
preserve streamline momentum. The exported products include:
\begin{itemize}
    \item \texttt{\_skeleton.tif}: traced binary skeleton;
    \item \texttt{\_pred\_edt.tif}: predicted localization confidence;
    \item \texttt{\_pred\_vis.tif}: predicted visibility;
    \item \texttt{\_pred\_orient\_conf.tif}: orientation-channel confidence.
\end{itemize}

\section{Validation Coverage}

The test suite focuses on the behavior most likely to regress:
\begin{itemize}
    \item STED target splitting between focus, annotation, and visibility bands;
    \item defocus and monomer/noise rendering behavior;
    \item real-data calibration profile construction and sampling;
    \item double-angle orientation handling;
    \item visibility-weighted losses and supervision floors.
\end{itemize}
Run the regression suite with:
\begin{quote}
\small
\begin{verbatim}
python -m unittest discover -s tests
\end{verbatim}
\end{quote}

\section{Typical Workflow}

A calibrated end-to-end run starts by measuring real images:
\begin{quote}
\small
\begin{verbatim}
python analyze_real_sted.py profile-real \
  --real_dir /path/to/real \
  --output_dir reports/sted_real \
  --patch_size 512
\end{verbatim}
\end{quote}
The resulting profile can drive synthetic data generation:
\begin{quote}
\small
\begin{verbatim}
python generate_dataset.py \
  --output_dir /path/to/sted2d \
  --bounds 1024 1024 \
  --synth_depth 16 \
  --calibration_profile reports/sted_real/sted_real_profile.json
\end{verbatim}
\end{quote}
Training then consumes the generated split directories:
\begin{quote}
\small
\begin{verbatim}
python train.py fit \
  --gpus 0 \
  --data_dir /path/to/sted2d \
  --dim 2 \
  --crop_size 512
\end{verbatim}
\end{quote}
A saved checkpoint can be evaluated with \texttt{train.py evaluate} or applied
to new images with \texttt{inference.py single}.

\end{document}
